\section{Lux FHE-C++: the production runtime}
\label{sec:fhe-cpp-production-runtime}

\texttt{github.com/luxfi/luxcpp/crypto/fhe} is the production runtime.
Its design priorities, in order, are: \emph{byte-equality with the
Go reference}, \emph{throughput on GPU}, \emph{a stable C ABI}, and
\emph{first-party-only code}.

\subsection{Surface}

The public C++ API lives under \texttt{cpp/include/lux\_fhe/}: opaque
parameter / key / ciphertext types, plus
$\mathsf{KeyGen}/\mathsf{Enc}/\mathsf{Dec}/\mathsf{Eval}$ entry points
parametrized by a \texttt{Backend} enum. The C ABI lives at
\texttt{c-abi/c\_lux\_fhe.cpp} and is the only surface Go cgo and
Rust bindgen consume. Downstream language bindings do not include C++
headers.

\subsection{Backends}

\begin{itemize}
\item \textbf{CPU (Lux FHE-Core CPU oracle).} Production code, byte-equal
to Lux FHE-Go. Lives under \texttt{cpp/backends/cpu/}. The CPU NTT, key-switch,
and bootstrap modules are the semantic oracles --- GPU backends are
byte-equal to them.
\item \textbf{Lux FHE-CUDA.} NVIDIA backend. Shipping (v0.1.0): the
NTT, key-switch, and programmable-bootstrap kernels live under
\texttt{cpp/backends/cuda/}. Backend selection is governed by the
compile-time \texttt{LUX\_FHE\_BACKEND\_CUDA} flag, the runtime
\texttt{cudaGetDeviceCount()} probe, and the
\texttt{LUX\_FHE\_BACKEND=\{cpu,cuda\}} environment override. The
device kernel output is asserted byte-equal to the CPU oracle on every
KAT entry by \texttt{fhe\_backend\_equiv\_test}.
\item \textbf{Lux FHE-HIP.} AMD backend. Mirrors Lux FHE-CUDA; lands
after the CUDA bootstrap stabilizes.
\item \textbf{Lux FHE-SYCL.} Portable accelerator backend. Reserved
name; ships only when a non-CUDA, non-HIP backend warrants it.
\end{itemize}

\subsection{Backend selection is explicit}

LP-167 \S\,Required claims for GPU-native credibility, point 9, forbids
silent dispatch. The C++ entry points take a \texttt{Backend} parameter
at handle creation; \texttt{LUX\_FHE\_AUTO} is intentionally undefined.
This is a discipline, not a UX failure: a runtime that silently falls
through from GPU to CPU loses the throughput contract that motivated
the GPU backend in the first place.

\begin{lstlisting}
#include "lux_fhe/keygen.hpp"

lux_fhe::SecretKey sk;
lux_fhe::PublicKey pk;
lux_fhe::EvaluationKey ek;
uint8_t seed[32]; /* derive from validator key material */

if (!lux_fhe::keygen_from_seed(lux_fhe::Backend::CUDA,
                               params, seed, sk, pk, ek)) {
    // CUDA unavailable -> caller decides whether to retry on CPU.
}
\end{lstlisting}

\subsection{Wire format}

The on-wire layout is fixed by \texttt{cpp/include/lux\_fhe/serialization.hpp}.
Magic / version / parameter-set-id / payload, little-endian. Lux FHE-Go
emits the same byte sequence; the cross-runtime KAT replay test
asserts byte-equality on each shipped parameter set. The wire format
is \emph{the} format --- there is no separate "C++ format" or
"Go format". Format parity is a normative claim
(LP-167 \S Required claims, point 1).

\subsection{What v0.1.0 ships}

The v0.1.0 milestone (this paper) covers:

\begin{itemize}
\item Public C++ API surface (every entry point callable);
\item C ABI surface for Go cgo / Rust bindgen;
\item CPU oracle for parameter resolution, deterministic keygen seed
expansion, NTT forward/inverse, gadget-decomposed key-switch, and
programmable bootstrap;
\item CUDA backend with byte-equal NTT, key-switch, and
programmable-bootstrap kernels; backend selected explicitly by the
compile-time flag \texttt{LUX\_FHE\_BACKEND\_CUDA}, the runtime
\texttt{cudaGetDeviceCount()} probe, and the environment override
\texttt{LUX\_FHE\_BACKEND}.
\item Encryption + decryption round-trip for $\mathsf{PN10QP27}$ and
$\mathsf{PN11QP54}$ on both backends;
\item Nebula-FHE transcript root (\texttt{lux\_fhe\_nebula\_root})
producing a deterministic 32-byte commitment;
\item KAT replay harness consuming the Go-generated corpus
(\texttt{fhe\_kat\_replay\_test}) and a cross-backend byte-equality
test (\texttt{fhe\_backend\_equiv\_test});
\item Microbenchmarks for NTT, key-switch, and bootstrap on CPU vs
CUDA dispatch (\texttt{luxcpp/crypto/fhe/bench/}). Reported numbers
in Section~\ref{sec:fhe-benchmarks}.
\end{itemize}

Lux FHE-Metal is scaffolded for v0.1.0 with the build wired up under
\texttt{LUX\_FHE\_BACKEND\_METAL}; production Metal device kernels reuse
the existing \texttt{luxcpp/crypto/ntt/gpu/metal/} kernel set and land
at v0.2.0.
